\documentclass[a4paper,12pt]{article}
\usepackage[T2A]{fontenc}
\usepackage[utf8]{inputenc}
\usepackage{cmap}
\usepackage[english, russian]{babel}
\usepackage{wrapfig}
\usepackage[backend=biber,style=gost-numeric]{biblatex} % стиль ГОСТ
\addbibresource{literatura.bib}
\usepackage{misccorr} % в заголовках появляется точка, но при ссылке на них ее нет
\usepackage{amssymb,amsfonts,amsmath,amsthm}  
\usepackage{indentfirst}
\usepackage[usenames,dvipsnames]{color} 
\usepackage[unicode,hidelinks]{hyperref}
% \hypersetup{%
%     pdfborder = {0 0 0}
\usepackage{multirow}
\usepackage{makecell,multirow} 
\usepackage{ulem}
%\usepackage{hyperref}
%\usepackage{graphicx,wrapfig}
%\renewcommand\epsilon{\varepsilon}
% \renewcommand\epsilon{\varepsilon}
%\graphicspath{{img/}}
\usepackage{geometry}
\geometry{left=1cm,right=1cm,top=2cm,bottom=2cm,bindingoffset=0cm,headheight=15pt}
\usepackage{fancyhdr} 
\linespread{1.2} 
\frenchspacing 
\renewcommand{\labelenumii}{\theenumii)} 
% \usepackage{caption}
%%%%%%%%%%%%%%%%%%%%%%%%%%%%%%%%%%%%%%%%%%%%%%%%%%%%%%%%%%%%%%%%%%%%%%%%%%%%%%%
%%%%%%%%%%%%%%%%%%%%%%%%%%%%%%%%%%%%%%%%%%%%%%%%%%%%%%%%%%%%%%%%%%%%%%%%%%%%%%%

\def\labauthor{Сарафанов И.Г.}
\def\labauthors{\labauthor}
\def\labnumber{127}
\def\labtheme{Таблица производных и интегралов}

%%%%%%%%%%%%%%%%%%%%%%%%%%%%%%%%%%%%%%%%%%%%%%%%%%%%%%%%%%%%%%%%%%%%%%%%%%%%%%%
	%применим колонтитул к стилю страницы
\pagestyle{fancy} 
	%очистим "шапку" страницы
\fancyhead{\LaTeX} 
	%слева сверху на четных и справа на нечетных
\fancyhead[L]{\labauthors} 
	%справа сверху на четных и слева на нечетных
\fancyhead[R]{Иродов \textnumero1.136} 
	%очистим "подвал" страницы
\fancyfoot{} 
	% номер страницы в нижнем колинтуле в центре
\fancyfoot[C]{\thepage} 
%%%%%%%%%%%%%%%%%%%%%%%%%%%%%%%%%%%%%%%%%%%%%%%%%%%%%%%%%%%%%%%%%%%%%%%%%%%%%%%
\usepackage{float}
\usepackage[mode=buildnew]{standalone}
\usepackage{tikz} 
% \usepackage{subcaption}
\usepackage{tikz,csvsimple,physics}
\makeatother
%\usepackage{booktabs}
\usepackage{pgfplots, pgfplotstable}
\usepackage[outline]{contour}
\usepackage{tocloft}
\renewcommand{\cftsecleader}{\cftdotfill{\cftdotsep}}


% вот этот удивительный мир

\begin{document}
\begin{center}
    \textbf{Домашняя работа \textnumero 3}
\end{center}
\section*{Задача}

\section*{Задача}

Два грузика массами $m/2$ соединены невесомым жёстким стержнем длины $l$.
Система покоится. В левый грузик налетает частица массы $m$
со скоростью $v$ и происходит абсолютно неупругий удар (частица
прилипает). Найти угловую скорость вращения системы после удара.

\section*{Решение}

После удара масса левого конца стержня становится
\[
m_1=m+\frac{m}{2}=\frac{3m}{2},
\]
масса правого конца
\[
m_2=\frac{m}{2}.
\]

Полная масса системы
\[
M=m_1+m_2=2m.
\]

\subsection*{1. Положение центра масс}

Пусть левый грузик имеет координату $x_1=0$, правый $x_2=l$.
Тогда координата центра масс равна
\[
x_{cm}=\frac{m_1x_1+m_2x_2}{m_1+m_2}.
\]

Подставляя значения, получаем
\[
x_{cm}=
\frac{\frac{3m}{2}\cdot0+\frac{m}{2}\cdot l}{2m}
=\frac{l}{4}.
\]

Следовательно, центр масс расположен на расстоянии $l/4$
от левого грузика.

\subsection*{Способ 1. Закон сохранения момента импульса}

Рассмотрим момент импульса относительно центра масс системы.

До удара движется только частица массы $m$,
поэтому её момент импульса относительно центра масс равен
\[
L_0=m v \frac{l}{4}.
\]

После удара система вращается вокруг центра масс:
\[
L=I\omega .
\]

Найдём момент инерции системы относительно центра масс.

Для левой массы
\[
r_1=\frac{l}{4}, \qquad m_1=\frac{3m}{2},
\]

для правой массы
\[
r_2=\frac{3l}{4}, \qquad m_2=\frac{m}{2}.
\]

Тогда
\[
I=
\frac{3m}{2}\left(\frac{l}{4}\right)^2
+
\frac{m}{2}\left(\frac{3l}{4}\right)^2.
\]

После вычисления получаем
\[
I=\frac{3ml^2}{8}.
\]

По закону сохранения момента импульса
\[
L_0=L.
\]

Следовательно,
\[
m v \frac{l}{4}
=
\frac{3ml^2}{8}\omega .
\]

Отсюда
\[
\boxed{\omega=\frac{2v}{3l}}.
\]

\subsection*{Способ 2. Кинематический подход с использованием скорости левого грузика}

Скорость центра масс найдём из закона сохранения импульса для всей системы. До удара импульс равен $mv$, после удара $M V_{cm}=2m V_{cm}$. Приравнивая, получаем
\[
V_{cm}=\frac{v}{2}.
\]

Рассмотрим теперь подсистему, состоящую только из налетевшей частицы и левого грузика, в момент удара. Внешние силы (реакция стержня) направлены вдоль стержня и не имеют горизонтальной составляющей, так как удар происходит перпендикулярно стержню. Следовательно, для этой подсистемы в горизонтальном направлении выполняется закон сохранения импульса:
\[
m v = \left(m + \frac{m}{2}\right) u_1,
\]
где $u_1$ — скорость левого грузика (с прилипшей частицей) сразу после удара. Отсюда
\[
u_1 = \frac{2v}{3}.
\]

После удара система движется как твёрдое тело: поступательно со скоростью центра масс $V_{cm}$ и вращается вокруг него с угловой скоростью $\omega$. Для левого грузика, находящегося на расстоянии $l/4$ от центра масс, его полная скорость складывается из скорости центра масс и линейной скорости вращения:
\[
u_1 = V_{cm} + \omega \cdot \frac{l}{4}.
\]

Подставляя найденные значения:
\[
\frac{2v}{3} = \frac{v}{2} + \omega \frac{l}{4},
\]
откуда
\[
\omega \frac{l}{4} = \frac{2v}{3} - \frac{v}{2} = \frac{v}{6},
\]
и окончательно
\[
\boxed{\omega = \frac{2v}{3l}}.
\]

Заметим, что правый грузик при этом оказывается мгновенно неподвижным, так как его вращательная скорость $\omega \cdot \frac{3l}{4} = \frac{v}{2}$ направлена противоположно $V_{cm}$ и полностью компенсирует поступательное движение.

\section*{Задача (абсолютно упругий удар)}

Два грузика массами $m/2$ соединены невесомым жёстким стержнем длины $l$.
Система покоится. В левый грузик налетает частица массы $m$
со скоростью $v$, направленной перпендикулярно стержню. Происходит
абсолютно упругий удар. Найти скорость частицы после удара,
скорость центра масс стержня и угловую скорость вращения стержня.

\section*{Решение}

После удара частица отскакивает со скоростью $u$ (положительное направление
совпадает с направлением $v$), стержень приобретает скорость центра масс
$V$ и угловую скорость $\omega$ (положительное направление -- против часовой
стрелки). Масса стержня с грузиками $M = 2m$, момент инерции относительно
центра масс (найден ранее) $I = \frac{3}{8}ml^2$, центр масс находится на
расстоянии $l/4$ от левого грузика.

Запишем законы сохранения.

\subsection*{1. Закон сохранения импульса}
\[
mv = mu + MV. \tag{1}
\]

\subsection*{2. Закон сохранения момента импульса}
Выберем неподвижную точку, совпадающую с положением левого грузика в момент
удара. До удара частица находится в этой точке, поэтому её момент равен нулю.
После удара частица всё ещё в той же точке, её момент также равен нулю.
Момент импульса стержня относительно этой точки складывается из момента
относительно центра масс и момента центра масс:
\[
L = I\omega + \frac{l}{4}\cdot M V,
\]
причём знаки выбраны так, что при $V>0$ (вверх) и $\omega>0$ (против часовой)
оба слагаемых положительны (векторное произведение даёт направление на нас).
Сохранение момента даёт
\[
I\omega + \frac{l}{4} M V = 0. \tag{2}
\]

\subsection*{3. Закон сохранения энергии}
\[
\frac{mv^2}{2} = \frac{mu^2}{2} + \frac{MV^2}{2} + \frac{I\omega^2}{2}. \tag{3}
\]

\subsection*{Решение системы}
Из (1) выразим $u = v - \frac{M}{m}V = v - 2V$.
Подставим в (3), предварительно умножив на 2:
\[
mv^2 = m(v-2V)^2 + MV^2 + I\omega^2.
\]
Раскрываем квадрат:
\[
mv^2 = m(v^2 - 4vV + 4V^2) + 2mV^2 + I\omega^2.
\]
Сокращаем $mv^2$:
\[
0 = -4mvV + 4mV^2 + 2mV^2 + I\omega^2,
\]
\[
4mvV = 6mV^2 + I\omega^2. \tag{4}
\]

Из (2) находим связь $\omega$ и $V$:
\[
I\omega = -\frac{l}{4}MV = -\frac{l}{4}\cdot 2m V = -\frac{ml}{2}V,
\]
откуда
\[
\omega = -\frac{ml}{2I}V = -\frac{ml}{2}\cdot\frac{8}{3ml^2}V = -\frac{4}{3l}V.
\]
Знак минус означает, что при $V>0$ угловая скорость направлена по часовой стрелке.
Для квадрата:
\[
\omega^2 = \frac{16}{9l^2}V^2.
\]

Подставляем в (4):
\[
4mvV = 6mV^2 + I\cdot\frac{16}{9l^2}V^2 = 6mV^2 + \frac{3ml^2}{8}\cdot\frac{16}{9l^2}V^2
= 6mV^2 + \frac{3}{8}\cdot\frac{16}{9}mV^2 = 6mV^2 + \frac{48}{72}mV^2 = 6mV^2 + \frac{2}{3}mV^2 = \frac{20}{3}mV^2.
\]
Сокращаем на $mV$ (полагая $V\neq0$):
\[
4v = \frac{20}{3}V \quad\Rightarrow\quad V = \frac{3}{5}v.
\]

Тогда
\[
\omega = -\frac{4}{3l}\cdot\frac{3}{5}v = -\frac{4}{5l}v,
\]
по модулю $\omega = \dfrac{4v}{5l}$.

Скорость частицы
\[
u = v - 2V = v - \frac{6}{5}v = -\frac{1}{5}v.
\]

\subsection*{Ответ}
\begin{itemize}
    \item Скорость частицы после удара: $\displaystyle u = \frac{v}{5}$,
    направлена противоположно начальной.
    \item Скорость центра масс стержня: $\displaystyle V = \frac{3v}{5}$,
    направлена так же, как начальная скорость частицы.
    \item Угловая скорость вращения стержня: $\displaystyle \omega = \frac{4v}{5l}$,
    направление вращения — по часовой стрелке (если смотреть сверху при ударе снизу).
\end{itemize}



















\end{document}